%% 05_discussion.tex

\section{Discussion}
\label{sec:discussion}

\paragraph{What the aggregate ordering actually measures.}
The A $>$ B $>$ C ordering on aggregate exact-match is real in the sense that it survives
a patient-level cluster bootstrap, but it is not a representation effect.  Under the
pre-specified scorer it is the product of a question bank in which 40.7\% of correct
answers are N/A and of a narrative that states explicitly which components a regimen
lacks; on questions with a substantive answer the three systems are within one point of
each other at 15--16\%, and every arm is below a constant per-family guess.  Under
final-answer scoring the narrative arm does lead on substantive questions, by 11--12~pp,
but that lead is confined to the families whose answers the narrative generation prompt
writes into the text and is absent on the one family both representations state equally.
A reader who wants a single sentence should take this one: the narrative arm's advantage
is made of explicit statements of absence, a summary that pre-computes answers, and a
scorer that treated the two kinds of response differently, and none of those is a
property of prose as a format.

\paragraph{Abstention is a first-class behaviour that benchmarks must stratify.}
The abstention result generalises beyond this dataset as a methodological point.  Any QA
benchmark in which ``not applicable'' is a valid answer, and in which templates fire
regardless of whether the queried entity exists, will have an N/A prevalence that varies
across strata (here, by tier).  Aggregate accuracy then mixes two skills, abstaining
correctly and answering correctly, with stratum-dependent weights, and a representation
that makes absence explicit will look better than one that does not.  We recommend that
clinical RAG evaluations report N/A prevalence per stratum, report accuracy separately on
N/A and substantive questions, and include a constant-N/A predictor alongside the usual
modal baseline.  Had we done this in the first version, the headline would have been
different.

\paragraph{Why the narrative arms abstain better.}
Both narrative generators produce a closed description of the regimen (``consists of one
medication component'').  The answer model, instructed to return N/A when the context does
not contain the answer, can act on an explicit negative.  Serialised FHIR resources are an
open-world representation: absence of a resource is not stated anywhere, and a question
about a missing component retrieves the nearest present one.  System~C's typed router
amplifies this by filtering the candidate pool to medication resources, so the context is
saturated with records of the component that does exist.  This is a genuine and, we think,
under-appreciated property of structured-record RAG: it needs an explicit mechanism for
answering ``there is no such thing'', for example a resource-count summary injected into
the context, or a pre-retrieval existence check.  The templated arm's stronger abstention
on the comparison family suggests that even the phrasing of the negative matters.

\paragraph{Summary leakage, not format.}
The generation prompt asks the narrative to list each component's first three and last
three administration dates, a mid-window date, the regimen start and end, and the total
administration count.  The substantive \texttt{temporal\_lookup} templates ask for the
first, last and next administration and the regimen start and end; the
\texttt{regimen\_aggregation} templates ask for counts; the \texttt{regimen\_compliance}
templates ask for windowed counts of the primary component.  These are the three families
on which the narrative arm leads under final-answer scoring, and the templated arm, which
writes the same fields in different words, leads on the same three.  On
\texttt{cross\_resource}, whose answers (component presence, schedule kind, tier) appear
verbatim in both representations, the arms tie.  The parsimonious explanation is that the
rendering step did part of the question answering.  The previous version attributed the
same pattern to ``FHIR JSON serialisation imposing a context-length cost that hurts
dose-count arithmetic''; we no longer think the evidence supports that reading over the
leakage one.  Separating them requires an aggregate-stripped narrative arm and a
structured arm restricted to the overlay resources in a compact rendering, which we list
as required follow-ups.

\paragraph{Where the structured arms hold their own.}
On \texttt{cross\_resource} substantive questions the structured arms match or edge the
narrative arm under both scorings, and System~C edges System~B on substantive questions
overall, consistent with its higher retrieval recall.  These are the questions for which a
single FHIR field is the answer.  The structured arms' aggregate deficit relative to the
narrative arm, and C's deficit relative to B, are N/A effects.

\paragraph{Retrieval versus reasoning at Tier~3.}
The earlier version located the Tier-3 problem in the answer LLM's reasoning, on the
grounds that retrieval recall stayed high while accuracy collapsed.  With the N/A
stratification the picture is simpler.  Substantive accuracy is low at every tier, and
Tier~3 is only somewhat lower than Tiers~1--2; what collapses at Tier~3 is the supply of
N/A questions.  Retrieval does contribute (about 13~pp over the no-retrieval arm at Tier~3
for every system under the pre-specified scorer), so ``actively misled by context'' was
the wrong description.  The more accurate statement is that a general-purpose LLM reading
either rendering of a multi-component regimen, without the anchor date it needs, under a
budget it frequently exhausts mid-reasoning, produces a correct windowed count or ratio a
small fraction of the time.

\paragraph{The three defects in perspective.}
The missing reference date (Defect~1) is a benchmark construction error that makes 7,800
questions unanswerable as posed and is the most consequential of the three; it must be
fixed by interpolating the date into the question before any rerun.  Thinking mode under
a 512-token cap (Defect~2) truncates roughly a third of responses before a final answer;
the 1024-token sweep shows this is not merely a budget problem, since 12--16\% of
responses are still truncated and substantive accuracy does not move.  A rerun should
either disable thinking (\texttt{reasoning\_effort=none}) or allow a budget of several
thousand tokens, and should state which.  The scorer's reach into the reasoning trace
(Defect~3) turned out to be far from neutral: it cost the narrative arm about 13~pp of
substantive accuracy and the structured arms about 2~pp, presumably because the
narrative's think blocks mention the total count and several dates before arriving at the
answer, so the ``first value'' heuristic picked the wrong one.  The free-text path's
failure on comparator and component-name answers is a plain bug that a corrected scorer
must fix.  We have chosen to publish the audit rather than delay until all three are
repaired, because the aggregate tables in the previous version have been cited and the
corrective is of value now.

\paragraph{Does this pairing hold anything constant?}
The honest answer is: the source, and the question set.  It does not hold information
content constant, because the narrative is a curated overlay summary with a precomputed
count and few dates, while the bundle is a complete record with every date and a large
irrelevant base.  Nor does it hold retrieval constant, because the narrative arm never
retrieves.  The comparison is therefore between an ``LLM summary in context'' pipeline and
two ``retrieve from a full record'' pipelines, which is a realistic deployment contrast
but not the causal isolation we originally described.  Disentangling representation from
summarisation requires at least three additional arms: a structured arm over the overlay
resources only, a compact structured rendering without JSON boilerplate, and a narrative
arm with aggregates and listed dates removed.

\paragraph{Why not a tool-calling arm?}
Ground truth is computed by deterministic functions over FHIR, so a SQL-on-FHIR or
function-calling arm would approach 100\% on this bank by construction and would not
inform the question of which text substrate a free-text RAG system should use.  The
free-text setting is motivated by ambient-scribe deployments in which narratives already
exist and case-manager questions arrive in natural language.  We accept that this
argument establishes relevance, not necessity: in a production PSP, a case manager's
count question should be answered by a query, not by an LLM doing arithmetic in its head,
and a tool-calling arm belongs in a follow-up as the reference ceiling against which any
retrieval pipeline's substantive accuracy should be judged.

\paragraph{Practical guidance.}
For practitioners the durable lessons are modest.  First, stratify by abstention before
believing any representation effect.  Second, if a structured record is the retrieval
substrate, add an explicit statement of what the record does not contain.  Third, if the
answer model is a hybrid reasoning model, decide deliberately whether thinking is on and
give it a budget that lets it finish; then score only the final answer.  Fourth, if a
summary is generated upstream of retrieval, recognise that whatever it pre-computes will
show up as ``QA accuracy'' downstream, and audit the summary prompt against the question
bank.  Fifth, a deterministic template is an acceptable narrative source for this kind of
regimen summary, which removes hallucination risk and generation cost.  Sixth, for
windowed adherence counts and ratios, compute them; do not ask a language model to.
